\section{ERP-Systeme in KMU}
\label{sec:erp_systeme_in_kmu}


ERP-Systeme (Enterprise-Resource-Planning-Systeme) stellen integrierte Anwendungssysteme dar, die unterschiedliche betriebliche 
Funktionsbereiche innerhalb einer gemeinsamen Systemarchitektur verbinden. Ziel ist die durchgängige Unterstützung und Koordination 
zentraler Unternehmensprozesse auf Basis einer einheitlichen Daten- und Informationsstruktur. Typischerweise umfasst dies Bereiche wie 
Beschaffung, Produktion, Lagerhaltung, Vertrieb, Finanzwesen oder Personalmanagement.

Im Unterschied zu isolierten Einzellösungen ermöglichen ERP-Systeme eine systemübergreifende Verknüpfung von Prozessen und Datenflüssen. 
Informationen werden nicht mehrfach in unterschiedlichen Anwendungen gepflegt, sondern innerhalb einer zentralen Datenbasis verarbeitet 
und verfügbar gemacht. Dadurch entstehen konsistente Informationsstrukturen sowie verbesserte Transparenz über betriebliche Abläufe und 
Abhängigkeiten.

Neben der operativen Prozessunterstützung übernehmen ERP-Systeme zunehmend auch eine infrastrukturelle Funktion innerhalb digitaler 
Unternehmensarchitekturen. Sie dienen häufig als zentrale Plattform zur Integration weiterer digitaler Anwendungen, beispielsweise im 
Bereich Reporting, Workflow-Automatisierung oder datenbasierter Entscheidungsunterstützung. Die Literatur beschreibt ERP-Systeme daher 
zunehmend nicht nur als betriebliche Standardsoftware, sondern als strukturelle Grundlage digital integrierter Organisationen.

Insbesondere in KMU erfolgt der Einsatz von ERP-Systemen häufig mit dem Ziel, organisatorische Komplexität zu reduzieren und Informationsflüsse 
besser zu koordinieren. Gleichzeitig zeigt die Forschung, dass die Einführung eines ERP-Systems nicht automatisch zu organisatorischer 
oder digitaler Transformation führt. Vielmehr hängt die tatsächliche Wirksamkeit maßgeblich davon ab, inwiefern das System mit bestehenden 
Prozessen, Strukturen und organisatorischen Voraussetzungen kompatibel ist.

